% =====================================================================
%  config/metadata.tex
%  ---------------------------------------------------------------
%  ÚNICO archivo que el equipo necesita editar para personalizar la
%  identificación del documento. No contiene contenido del plan.
% =====================================================================

% ---------------------------------------------------------------------
% 1. Identificación institucional (encabezado y pie de página)
% ---------------------------------------------------------------------
\newcommand{\uvUniversidad}{UNIVERSIDAD VERACRUZANA}
\newcommand{\uvFacultad}{Facultad de Estadística e Informática}
\newcommand{\uvPrograma}{Licenciatura en Ingeniería de Software}
\newcommand{\uvEE}{EE. Pruebas de Software}
\newcommand{\uvActividad}{PROYECTO FINAL}
\newcommand{\uvBandaHeader}{Proyecto de Pruebas}

% ---------------------------------------------------------------------
% 2. Identificación del documento
% ---------------------------------------------------------------------
\newcommand{\docTitulo}{Proyecto de pruebas de software}
\newcommand{\docProyecto}{\campo{Nombre del proyecto}}
\newcommand{\docFecha}{\campo{dd/mm/aaaa}}
\newcommand{\docEquipo}{\campo{Nombre del equipo}}
\newcommand{\docIntegrantes}{\campo{Matrícula y nombre de cada integrante}}
\newcommand{\docVersionPlantilla}{Ver. 01/10/25}   % versión de la PLANTILLA
\newcommand{\docVersionDoc}{\campo{0.1}}           % versión del DOCUMENTO del equipo

% ---------------------------------------------------------------------
% 3. Autoridad de aprobación
%    Quién autoriza desviaciones respecto de la estrategia acordada.
%    En el curso es el profesor; en un proyecto profesional cámbielo.
% ---------------------------------------------------------------------
\newcommand{\autoridadAprobacion}{\campo{Profesor de la Experiencia Educativa}}
\newcommand{\autoridadAceptacion}{\campo{Rol con autoridad para aceptar el riesgo residual}}

% ---------------------------------------------------------------------
% 4. Interruptores de la plantilla
%    Cambie a \...false para producir la versión final de entrega.
% ---------------------------------------------------------------------
\newif\ifmostrarinstrucciones
\mostrarinstruccionestrue      % \mostrarinstruccionesfalse  -> oculta las guías
\newif\ifmostrarejemplos
\mostrarejemplostrue           % \mostrarejemplosfalse       -> oculta los ejemplos

% «make final» compila una vez más definiendo \VERSIONFINAL, lo que apaga
% ambos interruptores sin tocar este archivo y deja el PDF de entrega en
% plan-pruebas-final.pdf. Editar las dos líneas de arriba produce el mismo
% resultado de forma permanente.
\ifdefined\VERSIONFINAL
  \mostrarinstruccionesfalse
  \mostrarejemplosfalse
\fi

% Las cajas de "Política académica" NO se ocultan con los interruptores
% anteriores: son contenido vigente del plan mientras el curso las exija.
% Si reutiliza la plantilla fuera del curso, elimine o sustituya esas cajas.
